% --- Archon physics formalization source begin ---
% archon:physics
% archon:covers IPhO2026Problems/problem_IPhO_2026_3_B_2.lean
% archon:source-report reports/ipho_2026_k3/problem_IPhO_2026_3_B_2.source.json
% archon:problem-id IPhO_2026_3
% archon:part-id B.2
% archon:previous-part IPhO_2026_3_A_3
% archon:previous-part-policy natural_language_prerequisite_only

\chapter{Physics problem IPhO\_2026\_3\_B\_2}
\label{ch:IPhO2026Problems_problem_IPhO_2026_3_B_2}

\paragraph{Problem source.}
Continue with the paramagnetic torus.  Its equation of state is T*M*V = n*K*H,
its heat capacity at constant M is C\_M = n*lambda/T\textasciicircum{}2, and dU = C\_M*dT.
The volume is fixed and the magnetic work on the material is
dW = mu\_0*V*H*dM.  Work and heat entering the torus are positive.

Current subquestion:
For an adiabatic change H\_i -> H\_f starting at T\_i, determine Delta T = T\_f - T\_i.

\paragraph{Current subquestion.}
For an adiabatic change H\_i -> H\_f starting at T\_i, determine Delta T = T\_f - T\_i.

\paragraph{Recorded answer/context.}
Delta T = T\_i*[sqrt((lambda + mu\_0*K*H\_f\textasciicircum{}2)/(lambda + mu\_0*K*H\_i\textasciicircum{}2)) - 1].

\paragraph{Figure/image path.}
/root/proposal\_for\_physic/science-mango/ipho\_2026\_source/image/T3\_page-3.png

\paragraph{Reusable previous-part conclusions.}
\begin{itemize}
\item Source A.3. Question: Subtract the vacuum-core contribution and write the work dW done on the paramagnetic material. Reusable conclusions: dW = mu\_0*V*H*dM. Policy: natural\_language\_prerequisite\_only; do\_not\_import\_Lean\_output
\end{itemize}

\paragraph{Formalization target.}
create a compiling Lean file with sorry bodies at `IPhO2026Problems/problem\_IPhO\_2026\_3\_B\_2.lean`.
The Lean declarations must preserve the physical quantities, dimensions or dimensional roles, figure labels, governing-law hypotheses, and final relation expressed by this problem.
Use Mathlib/Physlib names found through LeanExplore where available. If a domain API is missing, introduce faithful local abstractions rather than scalar placeholder aliases.

\begin{theorem}[Physics formalization target]
\label{thm:physics:IPhO_2026_3_B_2:target}
The assigned autoformalize agent should translate this physics problem into Lean declarations in the covered file, with theorem and lemma proof bodies written as `by sorry`.
\end{theorem}
\begin{proof}
This is an autoformalization task, not a proof task. Produce faithful statements that can later be proved without weakening the source contract.
\end{proof}
% --- Archon physics formalization source end ---
